\documentclass[12pt,a4paper]{article}

% Packages for math, code, and formatting
\usepackage{amsmath, amssymb, amsthm}
\usepackage{algorithm}
\usepackage{algpseudocode}
\usepackage{listings}
\usepackage{xcolor}

% Code style for C++
\lstdefinestyle{cpp}{
  language=C++,
  basicstyle=\ttfamily\small,
  keywordstyle=\color{blue},
  commentstyle=\color{green!50!black},
  stringstyle=\color{orange},
  numbers=left,
  numberstyle=\tiny,
  stepnumber=1,
  numbersep=5pt,
  frame=single,
  breaklines=true,
  tabsize=2,
  showstringspaces=false
}

% Theorem environments (for invariants, correctness proofs, etc.)
\newtheorem{theorem}{Teorema}
\newtheorem{lemma}{Lema}
\newtheorem{invariant}{Invariante}
\newtheorem{proposition}{Proposição}

\title{Relatório MC521 - Contest 3}
\author{Sofia Valverde Villas Bôas - 252974}

\begin{document}
\maketitle

\section*{G - Vapor Pressure}

Nesse problema, temos $A$ bolas e $B$ balões. Takahashi
remove repetidamente uma bola até que o número de
bolas seja no máximo $C$ vezes o número de balões. Queremos obter 
número de bolas ao final do processo, dividido pelo número de balões.

\subsection*{Solução}

Para a solução, basta remover bolas enquanto o número de bolas for maior
do que $C$ vezes o número de balões. Após isso, basta dividir o número de bolas
pelo número de balões e imprimir o resultado.

Depois, percebi que nem precisava ter feito o loop, já que, 
se $A$ for maior do que $C$ vezes $B$, o resultado é simplesmente $C$.
Caso contrário, o resultado seria simplesmente $A$ dividido por $B$.

Como o código já estava feito, deixei o loop mesmo.

\section*{E - Precondition }

Dadas duas strings $S$ e $T$, queremos determinar se $S$
satisfaz a condição: para cada letra maiúscula em $S$
que não seja a primeira, a letra imediatamente anterior a ela deve estar contida em $T$.

\subsection*{Solução}

A solução consiste em, inicialmente, colocar as letras de $T$ em um hash e,
depois, iterar sobre as letras de $S$ (exceto a primeira) e verificar se cada letra maiúscula
tem a letra anterior presente no hash.
Se encontrarmos uma letra maiúscula que não satisfaz a condição, podemos concluir que $S$ não satisfaz a condição
e imprimir "No". Caso contrário, se todas as letras maiúsculas satisfizerem a condição, imprimimos "Yes".

\end{document}